\cotefile{157-3}
\scanpage{162}
\mspage{111}
% archive sheet number: 162
\uline{Corollaire 2} (Moyennant \unsure{$W(7)$|$W\,7)$}) Soit $i:U\unsure{-|\to}X$ une imm. ouverte $W$-\unsure{dewyérienne|dwyérienne}. Pour que $i$ soit une $W$-cofibration triviale, il f. et s. que $i\in W$, ou ce qui revient au $=$, (\deleted{\unsure{prop.}} cf page 82) que $\dot T\hookrightarrow T$ soit dans $W$.

C'est évidemment nécessaire, et la suffisance résulte du th.~3 a).

\uline{Questions ouvertes} a) J'ignore si une immersion ouverte $i$ qui est une $W$-cofibration, voire une $W$-cofibration triviale, est forcément $W$-dewyérienne.\textsuperscript{*}
% slanted
\margin{\textsuperscript{*} et peut-être $\unsure{=}$ une imm. \unsure{Dwyer}!} Mais par \unsure{le|les} cor. de la prop 7 (p.~85), je sais que c'est le cas si on suppose $T$ isomorphe à la catégorie finale $e$. Il s'ensuit \deleted{\unsure{donc}} au moins ceci : Pour que non seulement $i:U-X$, mais aussi \deleted{\unsure{toute}} les immersions ouvertes $\unsure{i_1|u_1}:U-X_1$ qui s'en déduisent par le procédé \unsure{du|de} prop 7 b) au moyen de \deleted{\unsure{foncteurs}}\inserted{flèches} quelconques $T_1-T$ de but $T$ dans $\mathrm{Cat}$, soient des $W$-cofibrations, il f. et s. que \unsure{$i$|$u$} soit dewyérienne.

b) : \unsure{Y a-t-il} une bonne notion d'\uline{immersion \unsure{dwyérienne|dewyérienne}}, applicable à \uline{toutes} les immersions
\scanpage{163}
\mspage{112}
% archive sheet number: 163
et \unsure{incluant le|en dehors du} cas des imm. ouvertes et des imm. fermées ?

\uline{Démonstration du} th.~3 \deleted{\unsure{partie I}}\inserted{(p.~107)} Il suffit de prouver que si $i$ est \deleted{\unsure{$d$}}$W$-dewyérienne (\unsure{i.e.} $\unsure{i_0|i^*_0|i^*}$ \unsure{strictement} $W$-asphérique), alors $i'$ l'est également -- car il en résultera que $i'$ est $W$-asphérique et a fortiori $\in W$.

Il faut donc vérifier que pour \deleted{\unsure{$t$}} $t'\in\dot U'=\dot T$, l'élément $i'{}^*(t')$ \inserted{$=(u'\mapsto\operatorname{Hom}_{\unsure{X'|X}}(u',t'))$} de $U'{}^\wedge$ est \unsure{totalement} $W$-asphérique. Or $t'$ provient d'un unique élément $t$ de $\dot T$, et on sait qu'on a

\[
i'{}^*(\unsure{t|t'})=p_!(i^*(t)),\quad\text{où }p_!:U^\wedge\unsure{-|\to}U'{}^\wedge .
\]

Par hypothèse $F=i^*(t)$ est un élément tot. $W$-asphérique de $U^\wedge$, ce qui signifie que pour toute flèche $p:U-U'$ \inserted{dans $\mathrm{Cat}$} de source $U$, $p_!(F)$ est un élément $W$-asphérique de $\unsure{U'|U'{}^\wedge}$. En fait, on a le lemme suivant (par transitivité $(p'p)_!\simeq p'_!p_!$)

\uline{Lemme} \unsure{Soit}\deleted{\unsure{$F$}} $U$ dans $\mathrm{Cat}$, et \unsure{$F$|$R$} de $U^\wedge$. Si $F$ est tot. $W$-asph. \unsure{(resp.\textsuperscript{$W$}} ($U_{/F}$ est tot\textsuperscript{t} $W$-asphérique), alors pour tte flèche $p:U-U'$ dans $\mathrm{Cat}$ de source $U$, $F'=p_!(F)$ est un \unsure{objet} tot\textsuperscript{t} $W$-asph. de $U'{}^\wedge$ (resp.
\scanpage{164}
\mspage{113}
% archive sheet number: 164
$U'_{/F'}$ est un objet tot. $W$-asphérique de $\mathrm{Cat}$).

\uline{Corollaire 3} \inserted{(du th.~3)} Si $i$\inserted{$:U\unsure{-|\to}X$} est une immersion ouverte ou fermée $W$-dewyérienne (resp. \unsure{totalement} $W$-dewyérienne), alors pour toute flèche $p:U-U'$ dans $\mathrm{Cat}$, $i':U'\to X'=U'\amalg_{\unsure{U|U'}}\unsure{X'|X}$ l'est aussi. [\deleted{\unsure{C'est bien}}\inserted{NB} le cas non respé est essentiellement le théorème 3. Il \unsure{ne|en} dépend \unsure{d'aucun} axiome supplémentaire sur $W$, mais n'a d'intérêt qu'en vue des \unsure{équivalences} \textcircled{a} et \textcircled{b} (p.~108)\inserted{ qui en résultent}, et qui sont valables moyennant $W(7\mathrm{bis})$ resp. $W(7)$ seulement. ]

\uline{Commentaire} de \uline{motivation} : les immersions ouvertes ou fermées $W$-dewyériennes, ou \deleted{\unsure{les}} tot\textsuperscript{t} $W$-dewyériennes, \unsure{étant|sont} \deleted{\unsure{par th. st.}} stables par cochangement de base (d'après le corollaire précédent), mais n'ont pas du tout l'air de l'être par composition. A cet égard, les \lq\lq immersions de Dwyer\rq\rq{} (où on suppose l'existence de l'adjoint à droite, \deleted{\unsure{resp.}} \deleted{\unsure{à gauche}}, de $\unsure{v_0:|v_0\in} U-\overline U$, dans le cas d'une imm. ouverte) se comportent mieux, étant stables par composition. (\uline{NB} Il ne semble pas
\scanpage{165}
\mspage{114}
% archive sheet number: 165
non plus \inserted{clair} que les flèches str. $W$-asphériques, ou les tot. $W$-asphériques, soient stables par composition -- sorry !) Ma motivation pour m'échiner sur la notion $W$-dwyérienne est double. a) J'aimerais avoir des conditions \inserted{à la fois} \unsure{compréhensibles}, et aussi générales que possible, sur une immersion (ouverte p.~ex.), pour que ce soit une $W$-cofibration, \deleted{\unsure{notion}}\inserted{propriété} visiblement importante (tout comme la notion duale des $W$-fibrations).

b) J'ai envie, pour des applications ultérieures, de disposer d'une classe d'imm. ouvertes ayant les propriétés suivantes : 1º) Ce sont des $W$-cofibrations, et 2º) Si $U\overset{i}{\lhook\joinrel\unsure{-|\longrightarrow}}X$ est dans la classe, \deleted{alors} et si $\unsure{\varepsilon|E}\overset{\unsure{j}}{-}I$ est l'inclusion \inserted{\lq\lq sommets\rq\rq} de la catégorie discrète $\{0,1\}$ dans l'intervalle

\diagram{\begin{tikzcd}[ampersand replacement=\&,column sep=small]
\unsure{\mathscr{I}|\Phi}^{\deleted{\unsure{n|m}}}_{\unsure{n|u}}=\{0 \arrow[r] \& 1 \& 2 \arrow[l]\arrow[r] \& \cdots \& 2n-1 \& 2n \arrow[l]
\end{tikzcd}}
($j(0)\unsure{\mapsto|=}0,j(1)\unsure{=|\mapsto}2n$), \inserted{$n\unsure{\geq|>}1$,} l'inclusion \inserted{$\alpha$} de $(U\times I)\cup(X\times\varepsilon)$ dans $X\times I$ est dans la classe. Or cette \inserted{\unsure{amusante|double}} propriété \inserted{de stabilité} n'est \uline{pas} vérifiée par les immersions de Dwyer. (Prendre $i:\{0\}\longrightarrow\Delta^1=\{0,1\}$, l'inclusion $\alpha$ n'est pas \inserted{imm.} \unsure{Dwyer} \unsure{.|--}
\scanpage{166}
\mspage{115}
% archive sheet number: 166
2.5. Complément : \deleted{\unsure{intégration} d'un préfaisceau} \inserted{critère d'asphéricité totale} \deleted{sur}\inserted{d'}une catégorie \inserted{$W$-}propre sur une autre.

Soit

\[
f:X\longrightarrow I
\]
une catégorie propre sur une catégorie d'indices $I$. Si $I$ a un objet final $i_0$, alors \deleted{\unsure{le}} $X_{/i_0}$ s'identifie à $X$ et par suite l'inclusion \deleted{$X_{i_0}\longrightarrow X$ de}

\[
X_{i_0}\unsure{-|\longrightarrow}X
\]
de la fibre de $X$ en $i_0$ est $W$-asphérique, et \unsure{a|à} fortiori est dans $W$. Donc $X$ est asphérique \unsure{ssi|sss} $X_{i_0}$ l'est. Nous voulons \unsure{étendre|élever} ce résultat en un critère de \uline{totale} $W$-asphéricité, de façon précise prouver la

\uline{Proposition} 11 : Soit \deleted{\unsure{$I$|$f$}}\unsure{$f$}$:X-I$ $W$-propre, et \margin{$\searrow\!\searrow$}soit $i_0$ un objet final de $I$. Si $X_{i_0}$ est tot.\textsuperscript{t} $W$-asphérique, \inserted{(déf.~1, page \unsure{88|89})} il en est de même de $X$. (Sous réserve que $W$ satisfasse l'axiome de connexité $W(5)$.)

La démonstration sera donnée plus bas, après quelques préliminaires ayant un intérêt indépendant. Tout d'abord, rappelons que si $X$ est \inserted{$W$-}propre sur $I$, alors il donne naissance à un foncteur

\[
\begin{aligned}
I&\longrightarrow\mathrm{Hot}_W\\
i&\longmapsto\mathrm{hot}_W(X_i)\quad(\simeq\mathrm{hot}_W(X_{/i}))\,;
\end{aligned}
\]
en le composant avec le foncteur

\[
\pi_0:\mathrm{Hot}_W-\mathrm{Ens}
\]
(sous réserve que l'axiome \inserted{de connexité} $W(5)$, \deleted{\unsure{ait}} (p.~10) soit satisfait,
\scanpage{167}
\mspage{116}
% archive sheet number: 167
on trouve un foncteur

\[
\begin{aligned}
&\deleted{\unsure{\pi}}\\[-2pt]
i&\longmapsto\pi_0(X_i):\quad I\longrightarrow\mathrm{Ens}\\[-2pt]
&\hspace{3em}\tikz[baseline=-.4ex]{\node[rotate=90,inner sep=0pt] {$\simeq$};}\\[-2pt]
&\hspace{1em}\pi_0(X_{/i})
\end{aligned}
\]

Ceci posé, on a le résultat \unsure{subtil|crucial}

\uline{Proposition 12} : Soit $f:X\longrightarrow I$ $W$-propre, et supposons\textsuperscript{*} que $W$ satisfasse l'axiome de connexité $W(5)$. (sinon la $W$-propreté ne nous sera d'aucune utilité \ldots) On a alors une bijection canonique
% slanted
\margin{Il vaut mieux dire : Soit $f$ $W_0$-propre, où $W_0$ est le \unsure{localis\textsuperscript{t}} $0$-connexité \ldots{}}

\[
\pi_0(X)\xleftarrow{\sim}\underset{I}{\varinjlim}\,\pi_0(X_i),\tag*{\unsure{$(\alpha)$}}
\]
où la flèche est induite par les applications can. $\pi_0(X_i)\unsure{-|\to}\pi_0(X)$.

\uline{Dém} Passant à $f^\circ:X^\circ-I^\circ$, on est ramené à prouver le $=$ résultat pour $X$ \deleted{\unsure{lisse}}\uline{lisse} sur $I$, avec la seule différence que le foncteur

\[
i\longmapsto\pi_0(X_i)
\]
sur $I$ sera contravariant, i.e. un préfaisceau, ce qui offre l'avantage de l'intuition attachée au formalisme des préfaisceaux. Soit

\[
g:I\longrightarrow e,\qquad\text{d'où }\;gf=p:X-e,
\]
\deleted{\unsure{il}} est immédiat, par le calcul standard

\[
\deleted{\unsure{\pi}}\qquad p_!(F)=\pi_0(X_{/F})\qquad(F\in\mathrm{Ob}\,X^\wedge)
\]
\deleted{\unsure{que}} (on a identifié $e^\wedge$ à $(\mathrm{Ens})$), qu'on a

\[
\pi_0(X)=p_!(e_X)=g_!(f_!(e_X)).
\]

Or le préfaisceau

\[
i\longmapsto\pi_0(X_i)\quad(\simeq\pi_0(\unsure{{}_i\backslash X|i\backslash X}),\ \text{car }X\text{ \inserted{$W$-}lisse sur }I)
\]
\scanpage{168}
\mspage{117}
% archive sheet number: 168
n'est autre que $f_!(e_X)$. D'autre part, pour un préfaisceau quelconque \inserted{$G$} sur $I$, on a

\[
g_!(G)\simeq\pi_0(I_{/\unsure{C|G}})\simeq\underset{I}{\varinjlim}\,G(i),
\]
et appliquant ceci à $f_!(e_X)$, on trouve la formule \unsure{($\alpha$)} de la prop.~12.

Revenons à la situation \unsure{avec|où} $f$ $W$-\uline{propre}, et considérons un foncteur

\[
p:X\longrightarrow Y,
\]
nous sommes intéressés à \lq\lq calculer\rq\rq

\[
p_!(e_X)=(y\mapsto\pi_0({}_y\backslash X))\in\mathrm{Ob}\,Y^\wedge .
\]

Considérons pour ceci le diagramme suivant

\diagram{\begin{tikzcd}[ampersand replacement=\&]
{}_y\backslash X_i \arrow[r]\arrow[d] \& {}_y\backslash X \arrow[r]\arrow[d] \& {}_y\backslash Y \arrow[d]\\
X_i \arrow[r]\arrow[d] \& X \arrow[r,"\unsure{p|r}"]\arrow[d,"f"'] \& Y\\
\substack{\{i\}\\[-1pt]\simeq e} \arrow[r] \& I \& {,}
\end{tikzcd}}
où les trois carrés sont cartésiens. Comme $\,{}_y\backslash Y\unsure{-|\to}Y$ est $W$-propre, il en est de même de $\,{}_y\backslash X-X$ par chgt. de base, \unsure{donc} du composé $\,{}_y\backslash X-X\longrightarrow I$ par transitivité. Le diagramme montre d'autre part que la fibre de $\,{}_y\backslash X$ en $i\in\mathrm{Ob}\,I$ s'identifie à $\unsure{{}_y\backslash X_i}$.
\scanpage{169}
\mspage{118}
% archive sheet number: 169
La proposition 12 nous donne donc la formule

\[
\pi_0({}_y\backslash X)=\underset{I}{\unsure{\underline{\lim}|\varinjlim}}\,\pi_0({}_y\backslash X_i),
\]
laquelle formule ne peut s'empêcher d'être fonctorielle en $y$, \unsure{tout} comme celle de prop.~12 \deleted{\unsure{le}} l'est dans $\text{\uline{Prop}}_W\,I$. Or \deleted{\unsure{faisons}}\inserted{pour} $y$ variable dans $Y$, le \deleted{foncteur} préfaisceau sur $Y$

\[
y\longmapsto\pi_0({}_y\backslash X_i)
\]
n'est autre que $p_{i\,!}(e_{X_i})$, où

\[
p_i:X_i\longrightarrow Y
\]
est le composé $X_i\to X\xrightarrow{p}Y$ (induit par $p$).

On a ainsi prouvé

\uline{Corollaire} 1 (de prop.~12) Soit $f:X\to I$ $W$-propre,
% slanted
\margin{dire : $W_0$-propre \ldots} et $p:X\to Y$ un foncteur quelconque. Soit, pour tout $i\in\mathrm{Ob}\,I$, $p_i:X_i\to Y$ le foncteur induit sur la fibre \inserted{$X_i$} de $X$ en $i$. On a alors une loi \uline{fonctorielle} naturelle

\[
i\longmapsto p_{i\,!}(e_{X_i}):\quad I\longrightarrow Y^\wedge,
\]
et on a la formule

\[
p_!(e_X)\xleftarrow{\sim}\underset{I}{\varinjlim}\,p_{i\,!}(e_{X_i}),\tag*{($\alpha$)}
\]
où la flèche est induite par les morphismes canoniques

\[
p_{i\,!}(e_{X_i})\to p_!(e_X)
\]
\scanpage{170}
\mspage{119}
% archive sheet number: 170
induits par les inclusions $X_i\unsure{-|\to}X$.

La loi fonctorielle des $p_{i\,!}(e_{X_i})$ en $i$ a été explicitée en cours de démonstration. Notons que lorsque $X$ sur $I$ \deleted{\unsure{provient d'une}}\inserted{est cofibrée} et scindée, i.e. provient d'un foncteur

\[
i\longmapsto X_i,\quad I\deleted{\unsure{\to|\mathrm{Cat}}}\longrightarrow\mathrm{Cat},
\]
alors la loi fonctorielle de $p_{i\,!}(e_{X_i})$ en $i$ s'explicite de façon particulièrement \unsure{agréable} : à une flèche

\[
\alpha:i\to j,\quad\text{d'où }\alpha_*:X_i-X_j,
\]
est \unsure{associé|associée}

\[
p_!(\alpha):p_{i\,!}(e_{X_i})\longrightarrow p_{j\,!}(e_{X_j}),
\]
en notant qu'on a un diagramme

\diagram{\begin{tikzcd}[ampersand replacement=\&,column sep=small]
X_i \arrow[rr,"\alpha_*"] \arrow[dr,"p_i"',""{name=src,pos=.63}] \& \& X_j \arrow[dl,"p_j",""{name=tgt,pos=.63}]\\
\& Y \arrow[from=src,to=tgt,"\lambda_\alpha",bend left=7] \& {,}
\end{tikzcd}}
avec une flèche de commutation

\[
\lambda_\alpha:p_i\longrightarrow p_j\alpha_*
\]
(pas iso en \uline{général}), d'où

\[
p_{i\,!}(e_{X_i})\xrightarrow{(\lambda_\alpha)_!}
p_{j\,!}\underbrace{\alpha_{*\,!}(e_{X_i})}_{\longrightarrow e_{X_j}}
\longrightarrow p_{j\,!}(e_{X_j}),
\]
\unsure{plus} de doute c'est lui ! \margin{\tikz[baseline=-.2ex]{\draw[->] (0,0.65) .. controls (0,.22) and (.06,.13) .. (.52,0);} cette flèche devrait être explicitée avec soin !}
\scanpage{171}
\mspage{120}
% archive sheet number: 171
\uline{Corollaire 2} : Sous les conditions du cor.~1, supposons que $I$ ait un objet final $i_0$. Alors

\[
(\unsure{p_{i_0}|p\vert_{i_0}})_!(e_{X_{i_0}})\longrightarrow p_!(e_X)
\]
est un isomorphisme dans $Y^\wedge$.

D'où aussitôt la proposition 11 : si $X_{i_0}$ est \deleted{$Y_{/p_{i_0\,!}(e_{X_{i_0}})}$} tot. $W$-asph., alors \inserted{$X$ l'est, car} pour tout $p:X\to Y$, $p_!(e_X)\simeq p_{i_0\,!}(X_{i_0})$, et ce dernier élément de $Y^\wedge$ est $W$-asphérique par hyp. sur $X_{i_0}$.

Notons le cas particulier

\uline{Corollaire 3} Soit $Z$ \unsure{un|une} \deleted{catégorie}\inserted{objet} totalement $W$-asphérique dans $\mathrm{Cat}$. Alors son produit $Z\times I$ par \unsure{tte|tt} $I$ dans $\mathrm{Cat}$ ayant un objet final est aussi \inserted{tot.} $W$-asphérique.

\deleted{\uline{Corollaire \unsure{4}} Mais dans le cas d'une catégorie produit $X=Z\times I$, \unsure{considérée} comme cofibrée donc propre sur $I$, on peut faire mieux. Dans ce cas, le foncteur}

\[
\deleted{\unsure{i|j}\longmapsto p_{i\,!}(X_i)}
\]
\scanpage{172}
\mspage{121}
% archive sheet number: 172
3. \uline{Sommes amalgamées} \uline{et intégration}

\diagram{\(\displaystyle\int\!\begin{tikzcd}[ampersand replacement=\&,row sep=0pt,column sep=small]
\& X_1\\
X_0\arrow[ur]\arrow[dr] \&\\
\& X_2
\end{tikzcd}\,.\)}

\uline{3.1.} \uline{Réduction canonique au cas d'un diagramme d'immersions \deleted{\unsure{ouvertes}} fermées}.

Considérons un diagramme

\diagram{\(\text{(\unsure{$*$|$\alpha$})}\qquad\begin{tikzcd}[ampersand replacement=\&,row sep=small,column sep=large]
\& X_1\\
X_0\arrow[ur,"\unsure{\alpha_1|a_1}"]\arrow[dr,"\alpha_2"'] \&\\
\& X_2
\end{tikzcd}\)}
dans $\mathrm{Cat}$. Nous pouvons le regarder comme un contrafoncteur

\[
\unsure{\mathscr{I}|\Phi}^{\circ}\longrightarrow\mathrm{Cat}
\]
où

\diagram{\(\displaystyle\mathscr{I}=\left(\begin{tikzcd}[ampersand replacement=\&,row sep=0pt,column sep=small]
\& 1\arrow[dl]\\
0 \&\\
\& 2\arrow[ul]
\end{tikzcd}\right),\)}
donc comme définissant une catégorie fibrée (scindée) \inserted{$\unsure{\mathscr{X}|\mathcal X}$} sur $\mathscr{I}$. Il est parfois commode de représenter cette catégorie $\mathscr{X}$ par le \lq\lq signe d'intégration\rq\rq

\diagram{\(\displaystyle\int\!\begin{tikzcd}[ampersand replacement=\&,row sep=0pt,column sep=small]
\& X_1\\
X_0\arrow[ur]\arrow[dr] \&\\
\& X_2
\end{tikzcd}\,,\)}
et \unsure{la|le} regarder comme une sorte de \lq\lq somme amalgamée homotopique\rq\rq{} de $X_1$ et $X_2$ sur $X_0$, i.e. comme une \lq\lq limite inductive homotopique\rq\rq{} du diagramme de départ ($*$).
\scanpage{173}
\mspage{122}
% archive sheet number: 173
\unsure{Il} faut entendre par là que (pour $W$ \inserted{$\subset\mathrm{Fl}(\mathrm{Cat})$} donné) la classe \inserted{$\mathrm{hot}_W(\unsure{\mathscr{X}|\mathcal X})$} de $\mathscr{X}$ dans $\mathrm{Hot}_W$ a un \lq\lq sens homotopique\rq\rq{} remarquable \inserted{en termes de ($*$)} ; de façon précise, le diagramme ($*$) définit un objet $\unsure{\xi|\zeta}$ de $\mathrm{HOT}_W(\unsure{\mathscr{I}|\Phi})$ (qui est une catégorie de fractions de la catégorie de tous les diagrammes du type ($*$)\ldots), et si \deleted{\unsure{$\xi$}}

\[
\pi:\mathscr{I}\longrightarrow e
\]
désigne la projection absolue de $\mathscr{I}$, on a

\[
\pi_!(\xi)=\deleted{\unsure{h|H}}\mathrm{hot}_W(\mathscr{X}),
\]
d'après le calcul connu du foncteur $\pi_!$, adjoint à gauche du foncteur $\pi^*$

\[
\mathrm{HOT}_W(e)\overset{\text{déf}}{=}\mathrm{Hot}_W\longrightarrow\mathrm{HOT}_W(\mathscr{I}).
\]

% slanted
\margin{Cela suppose qu'on a $W(\unsure{3|2})$, non \unsure{simplem\textsuperscript{t}|seulement} $W(\unsure{2|3}\,\mathrm{bis})$.}

Or pour tout objet $I$ de $\mathrm{Cat}$, de projection absolue $p_I:I\to e$, le foncteur

\[
p_{I\,!}:\mathrm{HOT}_W(I)\unsure{-|\to}\mathrm{HOT}_W(e)=\mathrm{Hot}_W
\]
doit être regardé comme l'équivalent \lq\lq homotopique\rq\rq{} \unsure{,} du foncteur limite inductive $\underset{I}{\varinjlim}$ \unsure{des|de} diagrammes de type $I^\circ$ dans la catégorie $\mathrm{Hot}_W$ -- lesquelles limites inductives \uline{n'existent pas} dans $\mathrm{Hot}_W$ (sauf
\scanpage{174}
\mspage{123}
% archive sheet number: 174
pour des catégories très particulières -- les seules que je connaisse sont les catégories discrètes (qu'il faut peut-être $=$ supposer finies en plus).

Notons qu'on a un diagramme de foncteurs canoniques

\diagram{\(\unsure{(*)|(\alpha)|(**)}\qquad\begin{tikzcd}[ampersand replacement=\&,column sep=small,row sep=small]
\& X_0\arrow[dl,"\alpha_1"']\arrow[dr,"\alpha_2"] \&\\
X_1\arrow[dr,"{\substack{i_1\\ \deleted{\unsure{\alpha_2}}}}"'] \& \substack{\text{non}\\\text{commuta-}\\\text{tif !}} \& X_2\arrow[dl,"{\deleted{\unsure{\alpha_1}}\,i_2}"]\\
\& \unsure{\mathscr{X}|\mathcal X} \& {,}
\end{tikzcd}\)}
lequel diagramme n'a rien de commutatif dans $\mathrm{Cat}$. \uline{Il est cependant commutatif dans} $\mathrm{Hot}_W$ (quel que soit le localiseur fondamental $W$), \inserted{et même dans $\unsure{\mathrm{Cat}_{\mathrm{hot}}|\mathrm{Cat}\,\mathrm{hot}}$,} car si on désigne par

\[
i_0:X_0\longrightarrow\mathscr{X}
\]
l'inclusion canonique, on a des \unsure{hom|homs} canoniques

\diagram{\begin{tikzcd}[ampersand replacement=\&]
i_1\alpha_{\unsure{1|a}}\arrow[r] \& i_0 \& i_2\alpha_2\arrow[l]
\end{tikzcd}}
ce qui montre que $i_1\alpha_1$ et $i_2\alpha_2$ sont homotopes, donc définissent la $=$ flèche dans $\mathrm{Hot}_W$.
\scanpage{175}
\mspage{124}
% archive sheet number: 175
Considérons dans $\unsure{\mathscr{I}|\Phi}$ les deux \deleted{ouverts}\inserted{fermés} \inserted{isom. à $\Delta^1$}

\[
\begin{aligned}
\deleted{\unsure{U}}F_1&={}_1\backslash\mathscr{I}_{\unsure{/1|\deleted{/1}}}=\{0,1\}\\
\deleted{U}F_2&={}_2\backslash\mathscr{I}_{/2}=\{0,2\},
\end{aligned}
\]
on a

\[
\deleted{U}F_1\cap\deleted{U}F_2={}_0\backslash\mathscr{I}_{\deleted{\unsure{/0}}}=\{0\},\qquad
\deleted{U}F_1\cup\deleted{U}F_2=\mathscr{I}.
\]

Soient

\[
\begin{aligned}
\deleted{\unsure{\mathscr{X}|\mathcal X}_1={}}\,{}_1\backslash\mathscr{X}_{\deleted{/1}}&=\mathscr{X}\times_{\mathscr{I}}{}_1\backslash\mathscr{I}_{\deleted{/1}}=\mathscr{X}\times_{\mathscr{I}}\deleted{U}F_1\\
\deleted{\mathscr{X}_2={}}\,{}_2\backslash\mathscr{X}_{\deleted{/2}}&=\mathscr{X}\times_{\mathscr{I}}{}_2\backslash\mathscr{I}_{\deleted{/2}}=\mathscr{X}\times_{\mathscr{I}}\deleted{U}F_2
\end{aligned}
\]
les images inverses de $\deleted{U}F_1,\deleted{U}F_2$ dans $\mathscr{X}$, on a donc

\[
{}_1\backslash\mathscr{X}_{\deleted{/1}}\cap{}_2\backslash\mathscr{X}_{\deleted{/2}}=X_0,\qquad
{}_1\backslash\mathscr{X}_{\deleted{/1}}\cup{}_2\backslash\mathscr{X}_{\deleted{/2}}=\mathscr{X},
\]

\unsure{donc|d'où} le diagramme

\diagram{\(\text{($**$)}\qquad\begin{tikzcd}[ampersand replacement=\&,column sep=small]
\& \deleted{\unsure{\mathscr{X}}}\,X_0\arrow[dl,hook,"\alpha'_1"']\arrow[dr,hook,"\alpha'_2"] \&\\
{}_1\backslash\mathscr{X}_{\deleted{/1}}\arrow[dr,hook,"i'_1"'] \& \& {}_2\backslash\mathscr{X}_{\deleted{/2}}\arrow[dl,hook,"i'_2"]\\
\& \mathscr{X} \& {,}
\end{tikzcd}\)}
déduit du diagramme similaire \deleted{\unsure{des}}

\diagram{\begin{tikzcd}[ampersand replacement=\&,column sep=small]
\& \{0\}\arrow[dl,hook,no head]\arrow[dr,hook,no head] \&\\
\deleted{U}F_1\arrow[dr,hook] \& \& \deleted{U}F_2\arrow[dl,hook]\\
\& \mathscr{I} \& {,}
\end{tikzcd}}
est (comme ce dernier) un \deleted{diagramme}\inserted{carré} d'immersion\unsure{s} \deleted{\unsure{ouvertes|fermées}}\inserted{fermées} à la fois cartésien et cocartésien. D'autre part, on a un \unsure{hom}
\scanpage{176}
\mspage{125}
% archive sheet number: 176
canonique du diagramme \unsure{($*$)|($\alpha$)} dans ($**$), qui est l'identité \unsure{sur|pour} les sommets $X_0$ et $\unsure{\mathscr{X}|\mathcal X}$, et qui \deleted{provient du fait que} \unsure{sur|pour} $X_1,X_2$ est donné par les inclusions canoniques

\[
\begin{aligned}
X_1&\longrightarrow{}_1\backslash\mathscr{X}_{\deleted{/1}}\deleted{\unsure{{}_{/1}}},\qquad&
X_2&\longrightarrow{}_2\backslash\mathscr{X}_{\deleted{/2}}.\\
\Vert&&\Vert\\
\mathscr{X}_1\ \text{fibre}&&\mathscr{X}_2\ \text{fibre}\\
\text{de }\mathscr{X}\text{ en }1&&\text{de }\mathscr{X}\text{ en }2
\end{aligned}
\]

Ces inclusions, on \deleted{\unsure{le}} le sait, sont $\unsure{W|W_0}$-asphériques (elles ont $=$ un adjoint à droite), à fortiori elles sont dans $W$. Donc le diagramme non commutatif ($*$) est isomorphe, en tant que diagramme dans $\mathrm{Hot}_W$, au diagramme commutatif cartésien et cocartésien \inserted{($**$)} formé d'immersions \deleted{\unsure{ouvertes}}\inserted{fermées}. \deleted{\unsure{C'est|Ce}}Comme $\mathscr{X}$ est bel et bien somme amalgamée de ses deux \deleted{ouverts}\inserted{fermés} $\,{}_1\backslash\mathscr{X}_{\deleted{/1}},{}_2\backslash\mathscr{X}_{\deleted{/2}}$ sur leur intersection $X_0$, cela étaye l'intuition qui \deleted{\unsure{fait}} voit $\mathscr{X}$ comme une \lq\lq somme amalgamée \unsure{\lq\lq}homotopique\rq\rq. \unsure{$\S$}

Donnons \deleted{tout} de suite l'application suivante :

\uline{Proposition 1\deleted{\unsure{2|3}}\unsure{3|2}} : Soit un diagramme

\diagram{\begin{tikzcd}[ampersand replacement=\&,row sep=0pt,column sep=small]
\& X_1\\
X_0\arrow[ur,"\alpha_1"]\arrow[dr,"\alpha_2"'] \&\\
\& X_2
\end{tikzcd}}
dans $\mathrm{Cat}$, d'où le diagramme \unsure{($\alpha$)|($*$)} (p.~123)

\diagram{\begin{tikzcd}[ampersand replacement=\&,row sep=0pt,column sep=small]
\& X_1\arrow[dr,"i_1"] \&\\
X_0\arrow[ur,"\alpha_1"]\arrow[dr,"\alpha_2"'] \& \& \mathscr{X}\\
\& X_2\arrow[ur,"i_2"'] \&
\end{tikzcd}}
